\documentclass[14pt, a4paper]{article}
\usepackage{float}
\usepackage{pgf-pie}
\usepackage{pgfplots} 
\usepackage{geometry}
\usepackage{listings}
\usepackage{hyperref}
\usepackage{graphicx}
\usepackage{ragged2e}
\usepackage{color}
\usepackage{subfiles}
\usepackage{multicol}
\newgeometry{left=2.0cm, right=2.0cm, bottom=2.0cm, top=1.0cm}
% \usepackage{multirow}
\usepackage{fontspec}
% \setmainfont{Arial}
% \usepackage{authblk}
% \usepackage{orcidlink}


\renewcommand{\baselinestretch}{1.5}

% Special and customized commands
% Keywords macro
\providecommand{\keywords}[1]
{
  \small	
  \textbf{\textit{Keywords---}} #1
}


\title{xAIrus: A Novel Approach Revolutionizing SDLC by Embedding AI Throughout
the Entire Development Lifecycle}

\author{
    Iman Asgari Mosleh Abadi \and 
    Milad Farzaneh \and 
    Alireza Soltani Neshan
}

\begin{document}
\maketitle

\section{Abstract}

\section{Introduction}

Artificial Intelligence (AI) is transforming the software development landscape
by introducing innovative approaches to problem-solving, automation, and process
optimization. To systematically integrate AI into programming, this document
presents a dedicated Software Process Model and Methodology.

This methodology is designed to streamline and standardize the software
development process while boosting efficiency and clarity. By incorporating
clear criteria for each phase, it ensures that developers, teams, and
organizations can effectively utilize AI capabilities at every stage of the
software lifecycle, from inception to deployment.

With a structured, step-by-step explanation of the process model, this
methodology emphasizes adaptability, productivity, and quality. Each phase is
carefully crafted to enhance deliverability, set clear standards, and foster
collaboration, ensuring a transparent and efficient approach to software
development.

\tableofcontents

\section{The Software Process Model and Methodology}

The software process model and methodology for integrating AI into programming
comprises the following phases:

\subsection{Learn Requirements}

The "Learn Requirements" phase serves as the foundation of this methodology,
ensuring that the team has a deep understanding of essential concepts before
proceeding. This phase involves three critical components: learning about the
Software Development Life Cycle (SDLC), mastering Prompt Engineering, and
developing a thorough understanding of Large Language Models (LLMs). Each of
these components is crucial for effectively integrating AI into the software
development process.

\subsubsection{Learn About the Software Development Life Cycle (SDLC)}

The Software Development Life Cycle (SDLC) provides a structured approach to
software development. It is a framework that ensures the systematic progression
of a project through defined phases, leading to high-quality deliverables.
Understanding the SDLC is vital for developers and teams to establish a strong
foundation for AI-driven development.

\subsubsection*{Importance in AI Development}

By understanding the SDLC, teams can identify where and how AI tools and LLMs
can contribute, ensuring that AI enhances productivity without disrupting the
established process.

\subsubsection{Learn About Prompt Engineering}

Prompt Engineering is the practice of crafting effective and context-aware
prompts to maximize the output quality of LLMs. Since LLMs interpret prompts as
input to generate responses, mastering this skill is essential for leveraging AI
capabilities efficiently.

By mastering Prompt Engineering, teams can unlock the full potential of LLMs,
ensuring efficient and contextually accurate outputs across all phases of
development.

\subsubsection{Learn About Large Language Models (LLMs)}

Large Language Models (LLMs) are the backbone of AI-driven programming
methodologies. These models, trained on vast datasets, can generate human-like
text, assist in problem-solving, and provide insights across various domains.
Understanding LLMs' capabilities and limitations is key to effectively
incorporating them into the development lifecycle.

By gaining a deep understanding of LLMs, teams can assess their strengths and
limitations, ensuring that the chosen model aligns with the project’s
objectives.

\textbf{Outcome of the Phase:} Upon completing the "Learn Requirements" phase,
the team will have:

\begin{itemize}
    \item A solid grasp of the SDLC and its applicability in AI-driven development.
    \item Expertise in crafting effective prompts to maximize LLM efficiency.
    \item Comprehensive knowledge of LLMs’ capabilities, limitations, and practical
    applications.
\end{itemize}

This foundational knowledge will enable informed decision-making and lay the
groundwork for successful integration of AI into the software development
process.

\subsection{Choose LLM (Large Language Model)}

Selecting the appropriate Large Language Model (LLM) is a critical step in the
methodology, as it directly impacts the quality, efficiency, and capabilities of
AI integration in the software development process. The success of this step
relies on understanding the project requirements, evaluating available models,
and conducting thorough testing to ensure the chosen LLM aligns with the desired
outcomes.

\subsubsection{Evaluate Available LLMs}

The first step in choosing an LLM is evaluating the various models available in
the market. Each LLM comes with unique features, strengths, and limitations.
Consider the following factors during the evaluation:

\subsubsection*{Popular LLMs and Examples}

\begin{itemize}
    \item GPT-series models (e.g., GPT-4): General-purpose models suitable for
    natural language understanding, code generation, and content creation.
    \item Codex: Specialized in programming tasks, capable of generating code
    snippets, solving coding challenges, and debugging.
    \item Claude (Anthropic): Focuses on safety and reliability in natural
    language processing tasks.
\end{itemize}

\subsubsection*{Key Evaluation Criteria}

\begin{itemize}
    \item Model Size and Capabilities: Consider the model's size and training
    data. Larger models generally have better language understanding but may
    require more resources to operate.
    \item Training Dataset: Evaluate whether the model’s training data aligns
    with your project's domain. For example, Codex is trained extensively on
    public programming repositories, making it ideal for development tasks.
    \item Performance Benchmarks: Look at performance metrics for tasks like
    code generation, text summarization, or question answering, depending on
    your project needs.
    \item Cost and Licensing: Assess whether the model fits your budget, taking
    into account API usage, licensing fees, and deployment costs.
\end{itemize}

\subsubsection{Match the Model to Project Requirements}

After evaluating available LLMs, align their capabilities with the specific
requirements of your project. Consider these factors to ensure compatibility:

\begin{enumerate}
    \item Task Complexity:
    \begin{itemize}
        \item Simple tasks (e.g., generating documentation) can be handled by
        general-purpose models like GPT-3 or Claude.
        \item Advanced tasks (e.g., complex code generation or domain-specific
        optimizations) may require specialized models like Codex or a fine-tuned
        GPT-4.
    \end{itemize}
    \item Domain-Specific Needs:
    \begin{itemize}
        \item For projects requiring expertise in a specific domain, choose
        models fine-tuned for that purpose or consider fine-tuning a
        general-purpose model.
    \end{itemize}
    \item Deployment Environment:
    \begin{itemize}
        \item Determine whether the model will be used locally or accessed via
        an API.
        \item Local deployment may require open-source models like GPT-J or
        Llama 2 for better control over resources and customization.
    \end{itemize}
    \item Resource Constraints:
    \begin{itemize}
        \item Evaluate the computational resources needed to run the model.
        Larger models may require high-performance GPUs or cloud-based
        deployment.
    \end{itemize}
\end{enumerate}

\subsubsection{Test Shortlisted Models}

Once potential LLMs have been identified, conduct tests to verify their
suitability for the project. A structured testing phase ensures the chosen model
meets performance expectations.

\subsubsection*{Steps for Testing}

\begin{enumerate}
    \item Define Test Cases: Prepare a set of tasks representative of the
    project requirements, such as:
    \begin{itemize}
        \item Generating specific code snippets.
        \item Responding to technical queries.
        \item Summarizing project requirements.
    \end{itemize}
    \item Measure Performance: Evaluate the models based on:
    \begin{itemize}
        \item Accuracy: How well does the model generate correct and relevant
        outputs?
        \item Speed: How quickly does the model respond to inputs?
        \item Adaptability: Can the model handle variations in prompts
        effectively?
    \end{itemize}
    \item Compare Results: Rank the models based on their performance in the
    test cases and align the results with project goals. Tools for Testing:
    \begin{itemize}
        \item Use playgrounds provided by model providers (e.g., OpenAI, Hugging
        Face).
        \item Benchmark using APIs or local deployments with custom test
        scripts.
    \end{itemize}
\end{enumerate}

\subsubsection{Make the Final Selection}

Based on the evaluation and testing, select the model that best meets the
project’s needs. This involves balancing performance, cost, and compatibility.

\subsubsection*{Considerations for Final Selection}

\begin{itemize}
    \item Scalability: Ensure the model can handle growth in project complexity
    or usage.
    \item Support and Documentation: Choose models backed by robust community
    support and clear documentation for ease of integration.
    \item Customizability: Determine if the model supports fine-tuning or other
    customization options for better alignment with specific requirements.
\end{itemize}

\subsubsection*{Outcome of the Phase}

By the end of this phase, the team will have:

\begin{itemize}
    \item A detailed understanding of available LLMs and their capabilities.
    \item A tested and validated model that aligns with project requirements.
    \item Confidence in the selected LLM’s ability to enhance the software
    development process effectively.
\end{itemize}

This step ensures that the chosen LLM becomes a valuable asset, optimizing
workflows and enabling AI-driven innovation in the project.

\subsection{Develop Idea}

The "Develop Idea" phase is where the project's concept is refined, tailored,
and made actionable based on the team’s resources, skills, and potential. This
step takes into account whether the team already has a preliminary idea or
requires guidance to create one. Regardless of the starting point, this step
ensures that the project idea is localized, contextualized, and aligned with the
team's abilities and available resources.

\subsubsection{Input: Project Description and Resources}

To begin this phase, the following inputs are gathered:

\begin{itemize}
    \item Project Description: A high-level overview of the intended goals,
    objectives, and scope of the project.
    \item Team Skills: The expertise of team members in programming, AI, and
    relevant domains.
    \item Available Resources: Includes technical tools, computational power,
    budget, and access to data.
    \item Prompt Engineering Capabilities: The team's ability to design
    effective prompts for LLMs to ensure accurate and reliable results.
\end{itemize}

These inputs will drive the development of a practical, achievable, and
impactful project idea.

\subsubsection{Pathways for Idea Development}

\subsubsection*{Scenario 1: If the Team Already Has an Idea}

When a preliminary idea exists, this step focuses on refining it into a clear,
actionable concept. The refinement process involves:

\begin{itemize}
    \item Aligning the Idea with Team Skills: Evaluate the idea’s feasibility
    based on the team's expertise and available resources.
    \item Defining AI's Role: Specify how the LLM will be utilized, such as:
    \begin{itemize}
        \item Assisting in code generation.
        \item Providing decision support.
        \item Automating specific tasks (e.g., documentation or testing).
    \end{itemize}
    \item Identifying Gaps and Enhancements: Use prompt engineering techniques
    to clarify ambiguities in the idea and enhance it.
\end{itemize}

\textbf{Prompt Example:} \textit{"Based on the team's expertise in front-end
development and their access to GPT-4 API, refine the idea of building an
AI-powered customer support chatbot that integrates seamlessly into existing
e-commerce websites."}

\subsubsection*{Scenario 2: If the Team Does Not Have an Idea}

If the team lacks a concrete idea, this step involves generating recommendations
based on their potential and available resources. Prompt engineering can be
leveraged to solicit creative, relevant suggestions from an LLM.

\subsubsection*{Steps to generate an Idea}

\begin{enumerate}
    \item Assess the Team’s Potential: Identify strengths, such as:
    \begin{itemize}
        \item Familiarity with specific programming languages or frameworks.
        \item Experience in particular domains (e.g., healthcare, finance, or
        education).
    \end{itemize}
    \item Define Resource Constraints: Outline available tools, budget, and
    computational resources.
    \item Ask for Recommendations: Use well-crafted prompts to obtain AI-driven
    suggestions tailored to the team's context.
\end{enumerate}

\textbf{Prompt Example:} \textit{"Given a team proficient in Python, with access
to mid-level computational resources and a focus on developing educational
tools, suggest three project ideas where LLMs can be effectively integrated. The
ideas should consider the team's ability to create interactive and engaging user
experiences."}

\subsubsection{Output: Contextualized Project Concept}

The output of this phase is a Contextualized Project Concept, which encapsulates
the refined or newly generated idea tailored to the team’s context. This concept
should include:

\begin{itemize}
    \item Localized Relevance: How the idea aligns with the team's strengths,
    domain expertise, and resources.
    \item Feasibility: A preliminary assessment of whether the idea is realistic
    given the current resources and constraints.
\end{itemize}

\textbf{Outcome of the Phase:} By the end of this phase, the team will have:

\begin{itemize}
    \item A well-defined and actionable project concept tailored to their
    specific context.
    \item A solid foundation to transition into project planning and execution
    phases.
\end{itemize}

This step ensures that the project idea is not only innovative but also
practical, achievable, and optimized for success.

\subsection{Project Overview}

The "Project Overview" phase involves creating a high-level summary of the
project that aligns with the team's goals, skills, and available resources. This
step utilizes an iterative process, where input from the LLM is reviewed and
refined by the team to ensure alignment with the project’s requirements and
feasibility.

\subsubsection{Input: Contextualized Project Concept}

The input for this phase is the Contextualized Project Concept developed in the
previous step. This concept provides a foundation for creating a clear, concise,
and actionable project overview.

\subsubsection{Steps to Create a Project Overview}

\subsubsection*{Step 1: Get Project Overview from LLM}

The LLM is tasked with drafting a project overview based on the contextualized
project concept. This overview should include:

\begin{itemize}
    \item \textbf{Project Objectives:} The goals the project aims to achieve.
    \item \textbf{Major Deliverables:} Key outputs or features the project will
    produce. 
\end{itemize}

\textbf{Prompt Example:} \textit{"Based on the contextualized project concept
of} \texttt{[insert project concept]} \textit{, create a high-level project
overview. Include the project objectives, major deliverables, and the role of
the LLM in achieving these goals."}

\subsubsection*{Step 2: Analyze Project Overview in Team}

Once the LLM provides the draft overview, the team reviews it to ensure it meets
the following criteria:

\begin{enumerate}
    \item Alignment with Team Criteria: Does the overview reflect the team's
    skills, resources, and goals?
    \item Feasibility: Are the described objectives realistic given the
    available constraints?
    \item Clarity: Is the overview well-structured and easy to understand?
\end{enumerate}

\subsubsection*{Step 3: Refine with LLM if Necessary}

If the initial project overview does not meet the criteria, the team provides
feedback and refines the input prompts to guide the LLM more effectively. This
iterative process continues until the output aligns with the project’s
requirements.

\textbf{Refinement Prompt Example:} \textit{"The current project overview does
not fully reflect our team’s capabilities in front-end development or our
limited computational resources. Revise the overview to emphasize lightweight AI
integration and focus on user-friendly interface development."}

\subsubsection{Output: Validated Project Overview}

T he final output is a Validated Project Overview, which includes:

\begin{itemize}
    \item Justified Objectives: Clearly defined goals that reflect the team’s
    capabilities and the project’s intended impact.
    \item Specific Deliverables: A concise list of tangible outcomes expected
    from the project.
\end{itemize}

\textbf{Outcome of the Phase:} By the end of this phase, the team will have:

\begin{itemize}
    \item A polished and justified project overview that aligns with their
    contextualized project concept.
    \item A high-level summary ready to share with stakeholders, ensuring
    clarity and consensus across the team.
\end{itemize}

This step sets the stage for detailed planning by ensuring the project’s
foundation is solid, realistic, and well-communicated. The iterative loop
ensures that every aspect of the project overview is rigorously checked and
refined for optimal alignment.

\subsection{Tasks}

The "Tasks" phase involves breaking down the Validated Project Overview into
actionable and manageable tasks, formatted to align with the team's chosen
methodology (e.g., Agile, Waterfall, or other frameworks). This step ensures
that the work is organized and structured, providing clear direction to all team
members.

\subsubsection{Input: Validated Project Overview}

The input for this phase is the Validated Project Overview, which defines the
project’s objectives, deliverables, and AI integration details. This overview
serves as the foundation for creating tasks that reflect the project’s scope and
requirements.

\subsubsection{Steps to Create and Finalize Tasks}

\subsubsection*{Step 1: Get Tasks in the Desired Format from LLM }

The LLM is tasked with generating a list of tasks based on the validated project
overview. The output format should match the team’s preferred project management
methodology.

\subsubsection*{Examples of Desired Formats:}

\begin{itemize}
    \item Agile: Tasks categorized into epics, user stories, and sub-tasks.
    \item Traditional Waterfall: Tasks organized sequentially, corresponding to
    project phases (e.g., planning, design, development, testing).
    \item Custom Format: Any team-defined structure, such as prioritized task
    lists or milestone-driven deliverables.
\end{itemize}

\textbf{Prompt Example:} \textit{"Based on the validated project overview of}
\texttt{[insert project description]}\textit{, generate a detailed task list
formatted for Agile methodology.  Include epics, user stories, and sub-tasks for
each major deliverable. Ensure the tasks are specific, actionable, and aligned
with the project objectives."}

\subsubsection*{Step 2: Review Tasks Against Team Criteria}

Once the LLM generates the tasks, the team reviews the list to ensure it aligns
with:

\begin{enumerate}
    \item Relevance: Do the tasks address all aspects of the project overview?
    \item Clarity: Are the tasks clearly defined, specific, and actionable?
    \item Feasibility: Are the tasks realistic based on the team's skills and
    available resources?
    \item Prioritization: Are tasks properly prioritized based on the project's
    goals and dependencies?
\end{enumerate}

\subsubsection*{Step 3: Refine Tasks with LLM if Necessary}

If the initial list of tasks does not fully meet the team's needs, provide
feedback and refine the LLM's prompts to better guide the output. This iterative
process continues until the task list aligns with the project’s and team’s
criteria.

\textbf{Refinement Prompt Example:} \textit{"The generated tasks are too
high-level and lack actionable details. Refine the task list by breaking down
the epics into smaller, specific user stories and include acceptance criteria
for each task."}

\subsubsection*{Step 4: Finalize and Organize Tasks}

Once the tasks are fully refined, organize them into the chosen format. This may
involve:

\begin{itemize}
    \item Mapping tasks to a project timeline or sprint plan.
    \item Grouping tasks by priority, phase, or deliverable.
    \item Assigning ownership of tasks to specific team members or tools.
\end{itemize}

\subsubsection{Output: Actionable Task List (Desired Format)}

The output of this phase is an Actionable Task List, structured in the team’s
desired format. This task list should include:

\begin{itemize}
    \item Clear Objectives: Each task clearly supports the goals outlined in the
    project overview.
    \item Specific Details: Tasks are granular and actionable, with descriptions
    and acceptance criteria where applicable.
    \item Proper Organization: Tasks are grouped, prioritized, or sequenced
    based on the chosen methodology.
    \item Team Alignment: The task list reflects the team’s potential, ensuring
    realistic and achievable outcomes.
\end{itemize}

\textbf{Outcome of the Phase:} By the end of this phase, the team will have:

\begin{itemize}
    \item A detailed and structured task list that aligns with their project’s
    goals and methodology.
    \item A clear roadmap for executing the project, ensuring smooth transitions
    to subsequent phases.
    \item Flexibility to adapt tasks as needed while maintaining alignment with
    the overall project plan.
\end{itemize}

This phase ensures that the project’s workload is distributed efficiently,
enabling the team to work cohesively and achieve deliverables in a structured
manner. The iterative loop between LLM and team review ensures the tasks are
optimized for success.

\subsection{Development Roadmap}

The "Development Roadmap" phase involves organizing the Actionable Task List
into a structured and prioritized sequence of activities, providing a clear plan
for executing the project. The roadmap aligns tasks with the project’s
priorities and chosen methodology (e.g., Agile sprints, milestone-driven
development) to ensure smooth progression and efficient resource allocation.

\subsubsection{Input: Actionable Task List}

The input for this phase is the Actionable Task List generated in the previous
step. This list provides the foundation for creating a roadmap that sequences
and prioritizes tasks in alignment with the project's goals and constraints.

\subsubsection{Steps to Create and Refine the Development Roadmap}

\subsubsection*{Step 1: Get a Development Roadmap from the LLM}

The LLM is tasked with generating a development roadmap based on the actionable
task list. The format of the roadmap should align with the team's preferences,
such as:

\begin{itemize}
    \item \textbf{Sprint Plan:} Tasks grouped into sprints for Agile teams.
    \item \textbf{Milestone Plan:} Tasks sequenced based on major project
    milestones.
    \item \textbf{Comprehensive Prioritized List:} A single list of tasks
    organized by priority.
\end{itemize}

\textbf{Prompt Example:} \textit{"Based on the actionable task list, create a
development roadmap that prioritizes tasks for the first two sprints. Ensure
each sprint includes tasks of manageable scope, aligns with the project
objectives, and accounts for dependencies."}

\subsubsection*{Step 2: Adjust the Roadmap with the Team}

The team reviews the roadmap to ensure it meets the following criteria:

\begin{enumerate}
    \item Alignment with Priorities: Are the most critical tasks prioritized
    based on project requirements?
    \item Feasibility: Are the task groupings realistic given the team’s
    capacity and resources?
    \item Dependencies: Are task dependencies properly accounted for to avoid
    bottlenecks?
    \item Flexibility: Does the roadmap allow for adjustments based on evolving
    requirements or constraints?
\end{enumerate}

\subsubsection*{Step 3: Refine the Roadmap with LLM if Necessary}

If the roadmap does not fully align with the team’s needs, feedback is provided,
and the LLM is guided to revise it. This iterative loop continues until the
roadmap reflects the team’s priorities and constraints.

\textbf{Refinement Prompt Example:} \textit{"The initial roadmap does not
account for the team’s limited availability during the next sprint. Adjust the
roadmap to include fewer tasks per sprint while ensuring high-priority
deliverables are still met."}

\subsubsection*{Step 4: Finalize the Development Roadmap}

Once the roadmap meets the team’s criteria, finalize it by:

\begin{itemize}
    \item Organizing tasks into the selected format (e.g., Gantt chart, Kanban
    board, or sprint backlog).
    \item Communicating the roadmap to all stakeholders to ensure alignment and
    transparency.
\end{itemize}

\subsubsection{Output: Prioritized Tasks Based on Project Requirements}

The output of this phase is a Prioritized Development Roadmap, tailored to the
team’s needs and methodology. This roadmap should include:

\begin{itemize}
    \item Task Prioritization: Clearly defined order of tasks based on project
    goals and requirements.
    \item Grouping: Tasks organized into sprints, milestones, or phases for
    effective execution.
    \item Dependencies: Proper sequencing of tasks to account for
    interdependencies and resource availability.
    \item Flexibility: Built-in adaptability for handling unforeseen changes or
    adjustments.
\end{itemize}

\textbf{Outcome of the Phase:} By the end of this phase, the team will have:

\begin{itemize}
    \item A structured development roadmap that aligns with the project’s
    priorities and timeline.
    \item A clear sequence of activities, ensuring efficient use of resources
    and steady progress.
    \item Consensus among team members and stakeholders, minimizing confusion
    and delays.
\end{itemize}

This step ensures the project is set on a clear trajectory toward successful
delivery, with tasks prioritized and organized to support the team's workflow
and goals. The iterative loop between the LLM and the team ensures the roadmap
is optimized for the project’s specific requirements.

\subsection{Estimate Time and Milestones}

The "Estimate Time and Milestones" phase involves determining the time required
for each task and identifying key milestones to track project progress. This
phase provides a clear timeline for deliverables, ensuring the project stays on
schedule. It also integrates team feedback to refine the estimates and aligns
with methodologies like Agile by optionally mapping estimates to story points.

\subsubsection{Input: Tasks}

The input for this phase is the finalized Tasks, which provide a clear breakdown
of the work to be done. These tasks will be analyzed to determine time
estimates, milestones, and (if applicable) story points.

\subsubsection{Steps to Estimate Time and Milestones}

\subsubsection*{Step 1: Get Estimated Time and Milestones from LLM}

The LLM is tasked with providing an initial estimate for the time required to
complete each task and grouping tasks into milestones. If Agile is being used,
the LLM can also assign average story points to tasks for sprint planning.

\textbf{Prompt Example:} \textit{"Based on the following task list, estimate the
time required for each task in hours or days, and group tasks into milestones
based on logical deliverables. Additionally, assign story points to each task,
assuming an average team velocity of 20 story points per sprint."}

\subsubsection*{Step 2: Review Time and Milestones with the Team}

The team reviews the LLM’s time and milestone estimates to ensure they align
with:

\begin{enumerate}
    \item Team Skills and Capacity: Are the time estimates realistic given the
    team's skill level and availability?
    \item Task Complexity: Do the estimates appropriately reflect the complexity
    and dependencies of each task?
    \item Milestone Feasibility: Are the milestones logically grouped and
    achievable within the proposed timeline?
\end{enumerate}

\subsubsection*{Considerations for Story Points (if Agile):}

\begin{itemize}
    \item Evaluate whether the assigned story points are proportional to the
    task effort and complexity.
    \item Adjust the story points as needed based on the team's historical
    velocity and potential.
\end{itemize}

\subsubsection*{Step 3: Refine Estimates with LLM if Necessary}

If the initial estimates are not satisfactory, provide feedback to the LLM and
refine the estimates. This iterative process ensures the estimates are tailored
to the team’s potential and project needs.

\textbf{Refinement Prompt Example:} \textit{"The time estimates for tasks
involving front-end development seem too low given the complexity of the UI
features. Adjust the estimates for these tasks and ensure they align with the
team’s average completion rate of similar features in past projects."}

\subsubsection*{Step 4: Assign Tasks to Team Members}

Once the team is satisfied with the time estimates, milestones, and story points
(if applicable), assign tasks to team members. Assignments should:

\begin{itemize}
    \item Reflect each team member’s expertise and workload capacity.
    \item Balance the distribution of tasks to avoid bottlenecks.
    \item Align with sprint or milestone deadlines.
\end{itemize}

\subsubsection{Output: Task Time and Milestones}

The output of this phase is a detailed list of tasks with assigned:

\begin{itemize}
    \item Time Estimates: The time required to complete each task, expressed in
    hours, days, or weeks.
    \item Milestones: Logical groupings of tasks that serve as checkpoints for
    tracking project progress.
    \item (Optional) Story Points: If using Agile, story points are assigned to
    each task to aid in sprint planning.
    \item Task Assignments: Tasks are assigned to specific team members based on
    their expertise and availability.
\end{itemize}

\subsubsection*{Example of Output Structure}

\begin{enumerate}
    \item Task Name: Implement user authentication.
    \begin{itemize}
        \item Time Estimate: 16 hours.
        \item Milestone: Milestone 1 – Core Functionality.
        \item Story Points (if Agile): 8.
        \item Assigned To: Jane Doe.
    \end{itemize}
    \item Task Name: Create front-end dashboard.
    \begin{itemize}
        \item Time Estimate: 24 hours.
        \item Milestone: Milestone 2 – User Interface.
        \item Story Points (if Agile): 13.
        \item Assigned To: John Smith.
    \end{itemize}
\end{enumerate}

\textbf{Outcome of the Phase:} By the end of this phase, the team will have:

\begin{itemize}
    \item Realistic time estimates and well-defined milestones for tracking
    progress.
    \item A structured plan that aligns with team capacity and project goals.
    \item Balanced task assignments, ensuring all team members are contributing
    effectively.
\end{itemize}

This phase ensures that the project timeline is both realistic and achievable,
with clear milestones to measure progress and maintain alignment with the
project’s objectives. The iterative loop between the LLM and the team guarantees
the estimates are optimized for success.


\subsection{System Design}

The "System Design" phase focuses on creating a detailed design for the system
architecture, including any necessary diagrams and database schemas. This step
is essential for planning how various components of the system will interact,
ensuring scalability, efficiency, and alignment with project goals. The process
is collaborative and iterative, incorporating both team inputs and LLM-generated
suggestions to refine the design.

\subsubsection{Input: Finalized Project Overview}

The input for this phase is the Finalized Project Overview created in Step 4,
which provides a high-level summary of the project’s goals, deliverables, and
requirements. This serves as the foundation for extracting key design points and
building the system architecture.

\subsubsection{Steps to Develop the System Design}

\subsubsection*{Step 1: Extract Key Points of Design (Team Stage)}

\subsubsection*{Step 2: Get System Design from LLM}

\subsubsection*{Step 3: Review System Design (Team Stage)}

\subsubsection*{Optional Step 4: Get System Design Diagram from LLM}

\subsubsection*{Optional Step 5: Get Data Flow Diagram (DFD) from LLM}


\section{Conclusion}

The integration of Artificial Intelligence (AI) into the software development
process has the potential to revolutionize how teams approach programming,
collaboration, and delivery. This methodology provides a structured,
standardized, and efficient framework for leveraging AI at every stage of the
software lifecycle. By introducing clear criteria for each phase and adopting
best practices, this approach fosters clarity, productivity, and collaboration.

The outlined Software Process Model and Methodology aim to:

\begin{itemize}
    \item Boost the Development Process: Streamlining workflows and enhancing
    productivity through the strategic use of AI tools, such as prompt
    engineering and automated code generation.
    \item Ensure a Standardized Approach: Establishing a unified framework that
    promotes consistency, adaptability, and long-term maintainability across
    projects.
    \item Foster Clarity and Transparency: Providing well-defined roles,
    responsibilities, and criteria for each phase, ensuring alignment across
    teams and stakeholders.
\end{itemize}

By adhering to principles such as tiered review structures, reusable prompt
templates, and shared repositories for lifecycle documentation, this methodology
lays a strong foundation for AI-assisted software development. It not only
accelerates the development process but also sets clear standards for quality,
collaboration, and deliverability.

As software development evolves, this methodology serves as a guide to fully
harness the potential of AI, ensuring that teams can build innovative,
high-quality software in an efficient and transparent manner. This approach is
not just about incorporating AI but about transforming how we develop
software—making it faster, clearer, and more aligned with the demands of modern
technology.
\end{document}
